\section{Results}

% =======================================================================
\subsection{Dataset Statistics}
\label{section-dataset-diffs}
Before benchmarking any models on the collected dataset and Fin-SoMe, we assess how similar the datasets are using descriptive statistics. Figure \ref{figure-word-counts} shows the distribution of post lengths measured as the number of words per post. It indicates that tweets tend to be shorter than StockTwits messages on average. However, the distribution of their lengths is right-skewed suggesting that Twitter is occasionally used to share longer posts while StockTwits is not.


\begin{figure}[!ht]
	\includegraphics[width=\textwidth]{assets/images/word_count_distribution.pdf}
	\caption{Number of words per post for tweets (pyFin) and StockTwits posts (Fin-SoMe)}
	\label{figure-word-counts}
\end{figure}

The post length is not the only differentiating characteristic: The distribution of class labels also significantly differs between the two datasets. Within Fin-SoMe, more than 70\% of posts are bullish while only 10\% are bearish. While the data from Twitter also contains more positive than negative tweets, the class ratio is more balanced. However, this does not necessarily imply that posts on StockTwits are generally more positive. The Fin-SoMe dataset either uses slightly different label definitions or contains labeling errors. For example, the post \emph{``\$NXT.X December 28th is the key date. Dec 25-28 this is gonna be wild!''} was labeled as ``bullish''. However, no part of the sentence indicates a clear bullish sentiment. The post author could also expect a stock price drop. Given the ambiguous nature of the text, it would be labeled ``Uncertain'' according to the codebook used in this work.

\begin{figure}[!ht]
	\includegraphics[width=\textwidth]{assets/images/class_distributions.pdf}
	\caption{Distribution of labels within the datasets}
	\label{figure-class-distribution}
\end{figure}

Finally, the vocabulary used in the posts on the two platforms is not identical. The relative frequency of words like ``Elon'', ``chatroom'', ``inflation'', and ``factory'' is between 20 and 50 times higher in the Twitter dataset than it is in the posts from StockTwits. On the other hand, the words ``blackberry'', ``silver'', ``vix'', and ``senate'' are 12 to 25 times more prevalent in StockTwits posts than in tweets. The divergence between the two sets of vocabulary can be explained by platform-specific and temporal differences. Users on Twitter are more prone to interacting with one another, for example by commenting. Additionally, the Fin-SoMe dataset was collected in or before 2020 putting more than two years of time between the posts in both data sets. Therefore, the two datasets have different data distributions despite covering similar domains. These differences can affect a model's predictive performance, which will be assessed in more detail in Section \ref{section-modelperf-on-finsome}. The same concept applies to any model's ability to generalize to future, unseen data. If the nature of online discussions changes significantly, older models might require re-training on current data to maintain their performance characteristics.

% =======================================================================
\subsection{Model Performance Evaluation}

\subsubsection{Optimal Model Configuration}
All subsequently reported results will be based on the optimal model configuration that is found during hyperparameter optimization.
The optimal hyperparameter configuration for each model is presented in Table \ref{table-optimal-hparams}. The chosen hyperparameters suggest that for the ML models subword tokenization outperforms words-based tokenization using character n-grams of size 4. Another notable pattern is the DL models' need for high token dropouts. Both the recurrent neural net and the transformer-based neural net generalize best when 40 - 50\% of the input tokens are dropped. Additionally, the optimal embedding dimensionality is relatively low. A possible cause for this is that complex DL models require stronger regularization to not overfit smaller datasets.

\begin{table}[!ht]
\centering
\small

\begin{tabular}{lccccc}
\toprule
& \multicolumn{5}{c}{\textbf{Model}} \\
 \cmidrule(lr){2-6} 
\textbf{Hyperparameter} & Logistic Regression & SVM & Recurrent NN & Transformer NN & BERT Finetune \\
\midrule

vectorizer\_analyzer & char\_wb & char\_wb & -- & -- & -- \\
vectorizer\_ngram\_range & (4, 4) & (4, 4) & -- & -- & -- \\
vectorizer\_min\_df & 1.02e-06 & 2.87e-05 & -- & -- & -- \\
model\_C & 1.428 & 1.597 & -- & -- & -- \\
model\_kernel & -- & rbf & -- & -- & -- \\
model\_degree & -- & -- & -- & -- & -- \\
\midrule
token\_dropout & -- & -- & 0.498 & 0.439 & -- \\
embedding\_dim & -- & -- & 54 & 72 & -- \\
hidden\_dim & -- & -- & 189 & 237 & 238 \\
dropout & -- & -- & 0.414 & 0.286 & 0.249 \\
gru\_hidden\_dim & -- & -- & 44 & -- & -- \\
dim\_ff & -- & -- & -- & 140 & -- \\

\bottomrule
\end{tabular}
\caption{Optimal hyperparameter configuration for each model}
\label{table-optimal-hparams}
\end{table}




\subsubsection{Performance on Twitter Posts}
Figure \ref{figure-model-performance-twitter} displays each model's out-of-sample ROC AUC when applied to the dataset we collected from Twitter. The dictionary-based VADER and NTUSD-Fin perform worst scoring an AUC of 0.578 and 0.590 respectively. The domain-specific NTUSD-Fin outperforms VADER, but only by a small margin. FinBERT and TwitterRoBERTa, on the other hand, perform significantly better by both achieving an AUC score of 0.695.\newline
In comparison, all of the proposed models obtain AUC scores above 0.80. Out of these models, the recurrent neural net performs worst with an AUC of 0.803. It is outperformed by both the fine-tuned BERT model and the transformer-based neural network which score an AUC of 0.814. The two best models are the logistic regression and the SVM model which achieve AUC scores of 0.817 and 0.827 respectively.


\begin{figure}[!ht]
	\includegraphics[width=\textwidth]{assets/images/model_performance_pyfin.pdf}	
	\caption{Out-of-sample model performance on the dataset collected from Twitter}
	\label{figure-model-performance-twitter}
\end{figure}



\subsubsection{Performance on StockTwits Posts}
\label{section-modelperf-on-finsome}
To answer RQ4, we apply all models to the StockTwits posts compiled in the Fin-SoMe dataset. The corresponding model performances are dispalyed in Figure \ref{figure-model-performance-stocktwits}. The existing models' performance hardly changes compared to the ROC AUC on the tweet dataset. VADER now scores a slightly higher AUC of 0.590 while NTUSD-Fin achieves a score of 0.599. TwitterRoBERTa and FinBERT achieve AUC values of 0.705 and 0.710 respectively. On the other hand, all of the proposed models perform significantly worse on Fin-SoMe. The fine-tuned BERT model only scores an AUC of 0.702 while the transformer-based and recurrent neural network reach AUC scores of 0.707 and 0.709 respectively. The two ML models still perform best by obtaining an AUC of 0.727. Despite the relative performance decrease when applied to data from a different source, the proposed ML models still manage to outperform the existing BERT-based models. The performance decrease is likely due to the differences between the datasets as detailed in Section \ref{section-dataset-diffs}. The existing models are not affected by it because both datasets constitute out-of-distribution data for them.

\begin{figure}[!ht]
	\includegraphics[width=\textwidth]{assets/images/model_performance_finsome.pdf}	
	\caption{Out-of-sample model performance on Fin-SoMe}
	\label{figure-model-performance-stocktwits}
\end{figure}



\subsubsection{Inference Times}
In times where state-of-the-art language models like GPT-3 require millions of dollars in compute budget for training \cite{gpt3-cost}, it is important to assess the resourcefulness of machine learning models. Therefore, we quantify each models inference speed. Slower models require more compute time which, in a business setting, requires a larger compute budget. Figure \ref{figure-model-inference-times} presents the inference time per sample. Note that the y-axis is log-scaled due to the multiple orders of magnitude that lie between the fastest and the slowest model's speed. All experiments were conducted on a system with an AMD Ryzen 5 3600 CPU and 64GB of RAM.

\begin{figure}[!ht]
	\includegraphics[width=\textwidth]{assets/images/model_inference_time.pdf}	
	\caption{Inference time per sample (ms, log-scaled)}
	\label{figure-model-inference-times}
\end{figure}

VADER and the logistic regression model provide the fastest inference times below 0.1ms per sample. Both of the neural networks that we train from scratch are slower, but still deliver inference times of below 1ms, which is on-par with NTUSD-Fin. The SVM model is roughly 50 times slower than the fastest models at around 2.5ms per sample. Finally, all models that are based on LLMs are multiple orders of magnitude slower than the next quickest contender. At around 100ms per sample, they are more than 1000 times slower than the logistic regression. We have to acknowledge that all neural-network-based models would benefit from being run on a graphics processing unit (GPU). However, \shortciteA{buber2018performance} suggest that GPU utilization can speed up inference by a factor of 4 or 5, which would still leave the LLM-based models far behind all others in terms of speed.

\begin{itemize}[noitemsep]
	\item compare model performances on semEval?? financial phrase bank??
	\item mention training times?
\end{itemize}
% =======================================================================
\subsection{Model Diagnostics}
Considering the trade-off between training and inference speed an predictive performance, the logistic regression model 

\subsubsection{Topic-Specific Model Performance}
To gain a more in-depth understanding of the strengths and weaknesses of every model, we examine their performance of an array of several example tweets organized into four categories. During data labeling, these categories naturally emerged from the data. The most common topics that tweets can be categorized in are ``stock ownership'', ``options trading'', ``business acumen'', and ``neutral''. Posts subsumed under the ``stock ownership'' topic discuss investing decisions that involve buying or selling stocks directly as opposed to using more complex financial instruments. The ``options trading'' group refers to tweets discussing options trades such as buying or selling puts or calls. These posts tend to contain more complex language because stock options are a more sophisticated financial instrument than stocks. The ``business acumen'' category subsumes tweets that carry implicit sentiment that can only be detected by using business knowledge. Such tweets might discuss business events or decisions that imply a market sentiment, like exceeding earnings expectations or laying off staff. Finally, the category ``neutral'' serves as a control group because a good sentiment model should be able to detect documents without any sentiment and classify them as such. Table \ref{table-example-tweets} contains example tweets for each of the four topics and displays the correct label $\bm{y}$ as well as each model's predicted probability for the correct class $\bm{\mathbb{P}(\hat y = y)}$ which is emphasized if it was high enough to predict the label correctly.


\begin{table}[!ht]
\centering
\footnotesize
	\begin{tabular}{m{7.37cm}cccccc}
	\toprule
		\multirow{2}{*}{\textbf{Example Tweet}} & \multirow{2}{*}{$\bm{y}$} & \multicolumn{5}{c}{$\bm{\mathbb{P}(\hat y = y)}$} \\
		 \cmidrule(l){3-7}
		  &  & VADER & NTUSD & FinBERT & RoBERTa & pyFin \\
	\toprule


\emph{Stock Ownership} \\
\cmidrule(l){1-7}
adding to my \$AAPL position & POS & 0.00 & \emph{\textbf{0.52*}} & \emph{\textbf{0.66*}} & 0.21 & \emph{\textbf{0.91*}} \\
I'm long \$AAPL & POS & 0.00 & \emph{\textbf{1.00*}} & 0.05 & 0.44 & \emph{\textbf{0.91*}} \\
getting rid of my \$AAPL position & NEG & 0.00 & \emph{\textbf{0.56*}} & 0.05 & 0.16 & 0.36 \\
Just shorted \$AAPL & NEG & 0.00 & \emph{\textbf{0.65*}} & 0.04 & 0.24 & \emph{\textbf{0.97*}} \\
\cmidrule(l){1-1}
\emph{Options Trading} \\
\cmidrule(l){1-7}
Going all in \$TSLA 4/20 \$69 calls today before close & POS & 0.00 & \emph{\textbf{0.83*}} & 0.10 & 0.08 & \emph{\textbf{0.83*}} \\
Sold a 58P on \$INTC two weeks out & POS & 0.00 & \emph{\textbf{0.78*}} & 0.16 & 0.06 & 0.03 \\
Going all in \$TSLA 4/20 \$69 puts today before close & NEG & 0.00 & 0.46 & 0.03 & 0.01 & \emph{\textbf{0.93*}} \\
Sold a 58C on \$INTC two weeks out & NEG & 0.00 & 0.22 & 0.02 & 0.04 & 0.43 \\
\cmidrule(l){1-1}
\emph{Business Acumen} \\
\cmidrule(l){1-7}
\$TSLA factory can start production sooner than expected & POS & 0.00 & \emph{\textbf{0.75*}} & 0.11 & \emph{\textbf{0.81*}} & \emph{\textbf{0.75*}} \\
\$F beats EPS estimate, expected 1.34 reported 1.89 & POS & 0.00 & \emph{\textbf{0.93*}} & 0.07 & \emph{\textbf{0.53*}} & \emph{\textbf{0.63*}} \\
\$NFLX missed earnings estimates & NEG & 0.42 & 0.34 & \emph{\textbf{0.93*}} & \emph{\textbf{0.61*}} & 0.28 \\
\$OXY to lay off 42\% of staff leaked memo reveals & NEG & 0.20 & 0.10 & \emph{\textbf{0.96*}} & \emph{\textbf{0.49*}} & 0.29 \\
\cmidrule(l){1-1}
\emph{Neutral} \\
\cmidrule(l){1-7}
\$OXY stocks trading at \$123 & NEU & \emph{\textbf{1.00*}} & 0.00 & \emph{\textbf{0.80*}} & \emph{\textbf{0.87*}} & \emph{\textbf{0.77*}} \\
\$MMM Q3 numbers will decide the future of this stock & NEU & \emph{\textbf{1.00*}} & 0.00 & \emph{\textbf{0.93*}} & \emph{\textbf{0.86*}} & \emph{\textbf{0.54*}} \\
Come join our chatroom for exclusive stock tips!! \$CASH & NEU & \emph{\textbf{0.62*}} & 0.00 & \emph{\textbf{0.92*}} & 0.39 & \emph{\textbf{0.87*}} \\
Kathryn Janeway to take over as new \$SBUX CEO & NEU & \emph{\textbf{1.00*}} & 0.00 & \emph{\textbf{0.93*}} & \emph{\textbf{0.93*}} & 0.38 \\



	\bottomrule
	\end{tabular}
	\caption{The models predictions on example tweets organized by common topics. Numbers are predicted probabilities for the \emph{correct class}. \emph{\textbf{x.xx*}} indicates that the model predicted the correct class.}
	\label{table-example-tweets}
\end{table}

VADER fails to detect any sentiment in the documents but classifies all neutral tweets correctly. This is caused by its dictionary-based design where the model defaults to predicting neutral in the absence of any sentiment-laden words. The clearly sentiment-laden terms that occur in the documents (e.g. ``long'') are domain-specific, hence VADER can not pick them up. NTUSD-Fin's finance-specific dictionary makes it perform better, classifying all tweets in the ``stock ownership'' category correctly. However, it fails do detect the absence of sentiment and has a strong positivity bias which can be explained by the 4:1 ratio of bullish to bearish terms in the dictionary. FinBERT correctly identifies all neutral documents but rarely detects the correct sentiment within the other categories. RoBERTa performs exceptionally well in the ``business acumen'' category which is unsurprising as it has been trained on news headlines which these tweets closely resemble. It fails to classify the sentiment in the other two topics but accurately predicts the label for the neutral tweets. Finally, the proposed logistic regression model (``pyFin'') performs well across all categories. It struggles most with implicit sentiment and negation. For example, it misclassifies some of the ``business acumen'' tweets and does not detect that ``Sold a 58P'' constitutes positive market sentiment. As it is not able to account for interaction between words, it cannot comprehend that both ``sold'' and ``58P'' (a put option with a strike price of \$58) independently suggest negative sentiment but jointly imply a positive market sentiment. A more detailed error analysis of the proposed model is presented in Section \ref{section-model-error-analysis}.


\subsubsection{Learned Patterns}

To determine how small domain-specific models can outperform more complex generic models, we assess how the models adapted to the task of market sentiment classification. Exemplarily, we will scrutinize the proposed logistic regression model and the transformer-based neural network. For the logistic regression, we evaluate which tokens most strongly indicate the presence of each sentiment class. For the neural network, we visualize the word embeddings which capture semantic information regarding the sentiment class assignment for each word.

Table \ref{table-top-tokens-per-class} displays the tokens with the largest associated coefficients indicating these tokens are the most informative for predicting in favor of a specific class. Most of the tokens are domain-specific and do not carry sentiment in regular English language. However, within the domain of finance, they carry a clear sentiment. For example, the model learned that the words ``run'', ``buy'', or ``bull'' are indicative of positive sentiment. The same holds for the bearish class for which the tokens ``dump'', ``shor'', and ``red'' are indicative. The predictive tokens for the neutral class are not as coherent, but do contain several indicators for spam posts, for example ``chat'' which is often used for advertising paid chatrooms. Additionally, the table demonstrates the advantages of subword tokenization. A Twitter user could use the words ``buuullish'' or ``bulllish'' to emphasize a positive market sentiment. While word-based tokenization fails to detect that these are different versions of the word ``bullish'', the subword tokenization recognize that they refer to similar concepts due to the same prefix (``bull''), middle (``ulli''), and ending (``lish'').

\begin{table}[!ht]
	\centering

	\begin{tabular}{ll}
\toprule
\textbf{Class} & \textbf{Tokens Associated with Largest Coefficients}\\
\midrule
\textbf{Bullish} & run, buy,  rip,  cal, call, 999c, bull, ulli,  bul, llis, lish,  ath,  \includegraphics[height=1em]{assets/images/emoji_rocket.png} ,  up ,  buy \\
\textbf{Neutral} & tick, play, hart, name,  hit,  | ,  = , real,  ser, \includegraphics[height=0.9em]{assets/images/unicode_lightshade.png},  9:9,  sur, er? ,  or , chat, \\
\textbf{Bearish} &fall, eari, dump,  dum, rish, lowe,  dow, shor,  low,  red, hort, 999p, down,  los,  put, \\
\bottomrule
	\end{tabular}
	\caption{The top 15 tokens with largest model coefficients for each of the three classes}
	\label{table-top-tokens-per-class}
\end{table}

As opposed to the logistic regression, the neural networks we trained do not learn a single parameter per token. Instead, they learn word embeddings. For the transformer-based neural network, these word embeddings are 72-dimensional (see Table \ref{table-optimal-hparams}). Semantically, tokens that lie in close proximity to each other in the 72-dimensional embedding space are similar, i.e. associated with the same market sentiment label. To visualize the embeddings, the 72-dimensional vectors are projected into two a two-dimensional space which can be visualized using a scatterplot. For calculating this projection, we use t-distributed Stochastic Neighbor Embedding \shortcite{tsne-paper}, also known as t-SNE. It retains the local neighborhood structure in the data such that vectors that lie close together in the high-dimensional space also appear close together in the low-dimensional space which facilitates the identification of clusters. Figure \ref{figure-tsne-word-embeddings} shows this low-dimensional representation of selected word embeddings. The unit-less axes represent the two dimensions of the low-dimensional vector space.

\begin{figure}[!ht]
	\includegraphics[width=\textwidth]{assets/images/word_embeddings.pdf}
	\caption{Two-dimensional t-SNE representation of selected word embeddings}
	\label{figure-tsne-word-embeddings}
\end{figure}

The figure shows that the word embeddings manage to capture the semantics within and between the market sentiment classes. Words indicating positive market sentiment lie close together in the embedding space and the the different classes form distinct clusters.


\subsubsection{Model Error Analysis}

\label{section-model-error-analysis}
Look at errors and see that negation is the most prevalent cause!




% =======================================================================
\subsection{Analyzing Spam Posts(?)}

% -----------------------------------------------------------------------
% =======================================================================
\subsection{pyFin-Sentiment}
reference the python library.